% Source: https://artofproblemsolving.com/community/c3h3126203_power_of_a_point_and_uses
% Stable identity: aops-topic-3126203-q01
Parameterize a line from $P$ by signed distance $t$ along a unit vector $u$. Points on the circle satisfy $|P+tu-O|^2=r^2$, or $t^2+2u\cdot(P-O)t+PO^2-r^2=0$. Vieta makes the product of its two signed intersection distances $PO^2-r^2$. For an exterior secant directed toward the circle both roots are positive, giving the stated unsigned product. At tangency the radius is perpendicular to $PT$, so the Pythagorean theorem gives $PT^2=PO^2-r^2$. Hence in the example $PT=\sqrt{4\cdot9}=\boxed6$.
\begin{asy}
usepackage("amsmath");
size(7cm);
pair O=(0,0),P=(-5,0),A=(-3,0),B=(3,0),T=(-1.8,2.4);draw(circle(O,3));draw(P--B);draw(P--T--O,dashed);label("P",P,W);label("A",A,S);label("B",B,E);label("O",O,S);label("T",T,N);label("schematic",(0,-3.7));
\end{asy}
